\subsection*{MT-ART07. The entire non-real $A_5$ Artin category}
Let $L/\mathbb Q$ be Galois with group $A_5$, and suppose $L$ is not
totally real. Every nontrivial irreducible Artin L-function for this
extension is entire, unconditionally. Under $M_H$ each satisfies GRH.
This closes the non-real $A_5$ specialization of the holomorphy gate;
MT-ART05/06 remain useful for totally real extensions and other base fields.

The holomorphy assertion is classical, not a new finding. An exact primary
locator is Calegari, \emph{The Artin conjecture for some $S_5$-extensions},
\S1, publisher-author PDF p.~2, the paragraph immediately before
Theorem 1.1: it states the stronger automorphy assertion for $A_5$ when
complex conjugation is nontrivial, citing Khare--Wintenberger and Kim's
icosahedral functorial constructions. The two conditions are equivalent
here because a Galois extension has a real embedding exactly when every
embedding is real, equivalently complex conjugation acts trivially.
We cite that classical result rather than repeat its proof.

Here is the complete GRH transfer and its tensor-category form.
For every finite-dimensional representation $V$ of $A_5$, put
$m=\dim V^{A_5}$. Complete reducibility gives
\[
 L(s,V)=\zeta(s)^m
          \prod_{\chi\ne1}L(s,\chi)^{\langle V,\chi\rangle},
 \qquad \langle V,\chi\rangle\in\mathbb Z_{\ge0}.
\]
The product after removing $\zeta^m$ is entire by the cited theorem.
For each irreducible factor, Brauer induction expresses its L-function
as a finite product of integer powers of finite-order Hecke L-functions.
At a point of the open critical strip off its middle line, $M_H$ makes
every such factor finite and nonzero. The product is consequently nonzero
there, including when some Brauer exponents are negative. This proves
the claimed GRH without making a positivity assumption on those exponents.

At $s=1$ every nontrivial irreducible factor is also nonzero. Indeed in a
Brauer expression $\chi=\sum_j a_j\operatorname{Ind}_{H_j}^{A_5}\eta_j$,
the order there is $-\sum_{\eta_j=1}a_j$, by the simple Dedekind pole and
Hecke boundary nonvanishing. Frobenius reciprocity identifies the sum with
$\langle\chi,1\rangle=0$. Thus $L(s,V)$ has a pole of order exactly $m$
at $1$, and has no other poles. Tensor products, duals, exterior and
symmetric powers, and every Schur functor produce representations of the
same finite group, so the same conclusions hold for them after removing
their explicitly determined invariant summands. This is a statement about
the whole Artin tensor category attached to this cover, not about a
general positive-weight motive. The companions in Grothendieck's language
are SGA~1, Expos\'e~V, \S7 (finite monodromy), and SGA~6, Expos\'e~V,
\S\S2,7 (representation-ring operations), together with AG-LAMBDA.

\subsection*{MT-ART08. A finite exceptional set is enough for a holomorphy transfer}
Here is a useful precise comparison lemma. Let $V$ be a self-dual,
finite-image complex Galois representation over a number field $K$.
Suppose an Euler product $F(s)$ converges in a right half-plane and
continues to an entire function, with the same finite local factors as $L(s,V)$ outside a
finite set $S$. Suppose also that, for each $v\in S$, the reciprocal of
the local factor of $F$ is a polynomial $Q_v((Nv)^{-s})$.
Then $L(s,V)$ is entire. No equality of ramified local parameters is
needed for this conclusion.

To prove this, write the Artin factor at $v$ as
$P_v((Nv)^{-s})^{-1}$. Every eigenvalue of Frobenius on the inertia
invariants is a root of unity, since $V$ has finite image. Hence the zeros
of $P_v((Nv)^{-s})$ all lie on $\operatorname{Re}s=0$.
Equality in a right half-plane, followed by meromorphic continuation,
gives
\[
 L(s,V)=F(s)\prod_{v\in S}
             \frac{Q_v((Nv)^{-s})}{P_v((Nv)^{-s})}.
\]
The right side is holomorphic on $\operatorname{Re}s>0$.
The classical completed Artin functional equation is
\[
 A^{s/2}\Gamma_V(s)L(s,V)
   =W A^{(1-s)/2}\Gamma_V(1-s)L(1-s,V),
\]
with $A>0$ and $W\ne0$; the gamma factor is a product of
$\Gamma_{\mathbb R}(s)$, $\Gamma_{\mathbb R}(s+1)$ and
$\Gamma_{\mathbb C}(s)$ with nonnegative integral multiplicities.
An exact locator for these formulas is Prasad--Yogananda,
\emph{A Report on Artin's holomorphy conjecture}, \S3.3,
PDF pp.~8--9, with Brauer induction in Theorem 3.3.
For $\operatorname{Re}s<1$, the numerator
$\Gamma_V(1-s)L(1-s,V)$ is holomorphic, since
$\operatorname{Re}(1-s)>0$. The reciprocal $1/\Gamma_V(s)$ is entire.
The functional equation therefore makes $L(s,V)$ holomorphic on this
second half-plane too. The two half-planes cover $\mathbb C$, proving
the assertion. For a non-self-dual $V$ the identical argument works if
the hypotheses are also supplied for $V^\vee$.

The finiteness of $S$ is an actual hypothesis. Agreement at a set of
places of density one does not give the finite product in the proof.
In particular this lemma cannot by itself upgrade the weak
five-dimensional correspondence in Calegari's Theorem 1.1.
In geometric terms this separates modification at finitely many closed
points of an arithmetic base from a density statement about infinitely
many Frobenius fibres.

\subsection*{MT-S501 and MT-S502. The $S_5$ transfer and its one-value gate}
Let $L/\mathbb Q$ be an $S_5$ extension, with complex conjugation of cycle
type $(12)(34)$. Assume $5$ is unramified and its Frobenius also has this
cycle type. Then the irreducible factors of dimensions $4$ and $6$ are
entire and, under $M_H$, satisfy GRH (MT-S501).

Let $E=L^{A_5}$, let $F\subset L$ have degree $6$ over $\mathbb Q$, and
put $H=EF$. Under $M_H$, the additional single condition
\[
                         \zeta_H(1/2)\ne0
\]
implies that every nontrivial irreducible factor is entire and satisfies
GRH (MT-S502). This is a proved reduction with the displayed extra input,
not a proof of that central nonvanishing.

The classical inputs are Calegari's Theorems 1.1 and 1.2,
publisher-author PDF pp.~2--3; the definition $H=EF$ is explicit in
Theorem 1.2 of arXiv:1112.1152v1, p.~2. The former supplies the
dimension-$4$ and dimension-$6$ holomorphy, and the latter supplies the
remaining case if $\zeta_H$ has no zero in the real interval $(0,1)$.
Under $M_H$ the only possible real zero in that interval is $1/2$.
Thus its interval condition is exactly equivalent to the displayed
single-value condition. Once holomorphy is available, the Brauer
zero-location argument of MT-ART07 proves GRH for every factor.
This is the full proof of both transfers. The one-dimensional factors
are Dirichlet functions and already fall under $A_D$.

The totally real quadratic field $E$ is the base for the odd projective
two-dimensional representations used in the cited theorem. Accordingly
this is a higher-rank consequence of a specific automorphic descent
theorem. It is not obtained by treating every $S_5$ factor as a single
Hecke character. In the Artin-motive programme, the remaining gate in
this family is one central Dedekind value. AG-SIGN controls a nonzero
leading coefficient after a central zero is removed; it does not prove
that the central order is zero.

\subsection*{MT-S503. One unconditional holomorphy example becomes a full GRH example}
Take the splitting field $L$ of
\[
                        x^5-x^3-x^2+x+1.
\]
Calegari, Theorem 5.1, publisher-author PDF p.~15, proves automorphy of
every representation of $\operatorname{Gal}(L/\mathbb Q)=S_5$;
the quintic field has discriminant $1609$. Consequently under $M_H$
every nontrivial irreducible Artin L-function of this explicit field
satisfies GRH, by the same signed Brauer argument. This example uses
that cited theorem, not the local-at-$5$ hypotheses of MT-S501, which
the source explicitly says this polynomial does not satisfy.
The historical bounded-height computation in Theorem 5.3 is part of
the classical proof's evidence; it is not a computation performed in
this run. Geometrically this supplies a concrete finite \(S_5\)-cover
of the arithmetic base whose complete Artin category is covered.
